% \chapter{Практическая часть}
% \label{ch:chap2}

% \section{Match filter}

% \begin{lstlisting}
% int16_t* match_filter(int16_t* samples, int samples_count, int L, double* h, int* size){
%     *size = samples_count;
%     int16_t* new_samples = (int16_t*)malloc(samples_count * sizeof(int16_t)); 

%     //convolve
%     for(int n = 0; n < samples_count; ++n){
%         int16_t tmp = 0;
%         for(int m = 0; m < L; ++m){
%             if (n - m >= 0){
%                 tmp += samples[n-m]*h[m];
%             }
%         }

%         new_samples[n] = tmp;
      
%     }

%     return new_samples;
% }

% \end{lstlisting}

% Функция принимает массив семплов, размер этого массива, число L, которое показывает, сколько семплов приходится на бит, ИХ h, переменную
% size, в которой вернется размер массива на выходе. Функция выполняет дискретную свертку входного сигнала с ИХ.

% \section{Downsampling}

% \begin{lstlisting}
% int16_t* down_sampling(int16_t* samples, int samples_count, int L, int start){
%     int symbols_size = samples_count / L - start + 1;

%     int16_t* symbols = (int16_t*)malloc(symbols_size * sizeof(int16_t));
    
%     for(int i = start; i < symbols_size; ++i){
%         symbols[i] = samples[i*L];
%     }

%     return symbols;

% }
% \end{lstlisting}

% Функция принимает массив исходных семплов, размер массива, число L, которое указывает на то, с каким периодом мы будем делать 
% downsampling (сохраняем только каждое L-ое значение). Параметр start - индекс, с которого функция начнет брать каждый L-ый элемент.
% В функции итерируемся по новому массиву и заполняем его каждым L-ым значением.

% \section{downscale}

% \begin{lstlisting}
% double max(int16_t* arr, int arr_size){
%     double max = -__DBL_MAX__;
%     for(int i = 0; i < arr_size; ++i){
%         if(arr[i] > max){
%             max = arr[i];
%         }
%     }

%     return max;
% }

% double* downscale(int16_t* samples, int samples_count){
%     double* down_scale_samples = (double*)malloc(samples_count * sizeof(double));

%     double max_value = max(samples, samples_count);

%     for(int i = 0; i < samples_count; i+=2){
%         down_scale_samples[i] = samples[i] / max_value;
%     }

%     return down_scale_samples;
% }
% \end{lstlisting}

% Для того, чтобы дальше работать с символами, нужно нормализовать их, т.е привести к значениям, которые будут близки к значениям из
% отрезка [-1;1]. В качестве нормирующего коэффициента будет использоваться символ с наибольшим значением. Функция max используется
% для поиска символа с максимальным значением. В функции downscale итерируемся по исходному массиву символов и делим каждый на максимальное
% значение среди символов.

% \section{Quantization}

% \begin{lstlisting}
% int16_t* BPSK_quantizater(double* symbols, int symbols_size){
%     int16_t* rounded_symbols = (int16_t*)malloc(symbols_size * sizeof(int16_t));

%     int point0[] = {-1, 0};
%     int point1[] = {1, 0};

%     double distance0 = 0, distance1 = 0;

%     double dx = 0,dy = 0;

%     for(int i = 0; i < symbols_size; i+=2){
%         dx = point0[0]-symbols[i];
%         dy = point0[1]-symbols[i+1];

%         distance0 = sqrt(dx*dx + dy*dy);

%         dx = point1[0]-symbols[i];
%         dy = point1[1]-symbols[i+1];

%         distance1 = sqrt(dx*dx + dy*dy);

%         if(distance0 > distance1){
%             rounded_symbols[i] = 1;
%             rounded_symbols[i+1] = 0;
%         }else{
%             rounded_symbols[i] = -1;
%             rounded_symbols[i+1] = 0;
%         }
%     }

%     return rounded_symbols;
    
% }
% \end{lstlisting}

% Функция принимает массив искаженных символов и размер массива. В функции в виде массивов размером в 2 элемента задаются истинные точки,
% которые есть в BPSK созвездии. Далее в цикле перебираем символы (точки на созвездии) и вычисляем расстояние между искаженной точкой
% и заданными истинными точками. В результирующий массив запишем координаты истинной точки, с которой искаженная точка имеет наименьшее расстояние.

% \section{Demapper}

% \begin{lstlisting}
% int8_t* BPSK_demodulator(int16_t* symbols, int symbols_size){
%     int bits_size = symbols_size/2;

%     int8_t* bits = (int8_t*)malloc(bits_size * sizeof(int8_t));

%     int j = 0;

%     for(int i = 0; i < symbols_size; i+=2){
%         if(symbols[i] == 1){
%             bits[j++] = 1;
%         } else{
%             bits[j++] = 0;
%         }
%     }

%     return bits;
% }
% \end{lstlisting}

% Функция принимает массив округленных символов и размер массива. Функция итерируется по массиву и с помощью условий выбирает,
% какой бит сопоставить символу.

% \section{Test}

% Для проверки корректности работы приемника сгенерируем семплы на стороне передатчика и запишем в файл, а на приемной стороне
% считаем файл и попытаемся превратить семплы в информацию, которую отправлял передатчик. Попытаемся передать символ "H"

% \begin{lstlisting}
% int main(){
    
%     size_t samples_count = 0;

%     int16_t* samples = read_pcm("bpsk_samples.pcm", &samples_count);

%     if(samples == nullptr){
%         printf("File not found!\n");
%         exit(1);
%     }

%     printf("samples:\n");

%     for(int i = 0; i < samples_count; ++i){
%         printf("%d ", samples[i]);
%     }

%     printf("\n");

%     int h[] = {1,1,1,1,1,1,1,1,1,1};

%     int conv_size = 0;
%     int* conv = match_filter(samples, samples_count, 10, h, &conv_size);

%     printf("After convolve:\n");

%     for(int i = 0; i < conv_size; ++i){
%         printf("%d ", conv[i]);
%     }

%     printf("\n");


%     int* ds_samples = down_sampling(conv, samples_count, 10, 0);

%     printf("Samples after down sampling:\n");

%     for(int i = 0; i < conv_size/10; ++i){
%         printf("%d ", ds_samples[i]);
%     }

%     printf("\n");

%     double* down_scale_samples = down_scaler(ds_samples, samples_count/10);

%     printf("Samples after downscaling:\n");

%     for(int i = 0; i < conv_size/10; ++i){
%         printf("%f ", down_scale_samples[i]);
%     }

%     printf("\n");


%     int16_t* q_samples = BPSK_quantizater(down_scale_samples, samples_count/10);

%     printf("rounded symbols:\n");

%     for(int i = 0; i < samples_count/10; ++i){
%         printf("%d ", q_samples[i]);
%     }

%     printf("\n");

%     int8_t* bits1 = BPSK_demodulator(q_samples, samples_count/10);

%     printf("bits:\n");

%     for(int i = 0; i < samples_count/20; ++i){
%         printf("%d ", bits1[i]);
%     }

%     printf("\n");

%     int str_len = 0;
%     char* str = btos(bits1, samples_count/20, &str_len);

%     printf("string: %s\n", str);

%     return 0;
% }
% \end{lstlisting}

% \begin{figure}[H]
%     \centering
%     \includegraphics[width=1.0\textwidth]{rx_string.png}
%     \caption{Процесс обработки принятого сигнала}
% \end{figure}

% Приемнику удалось верно преобразовать семплы в текст, значит, удалось верно реализовать логику работы приемника. \\

% \section{Сигнал на выходе match filter}


% \section{Глазковая диаграмма}



% \endinput
